
\subsection{The Intermediate Argument} 
An essential ingredient of our proof is an 'intermediate argument', stating that prior to having a large slice that closes a non-trivial loop on the toric, there must be a step in which there is a slice that on one hand doesn't close a non-trivial loop, yet it has a size that is comparable to the toric side. Before the first step of that kind, we can assume that no slice closes a non-trivial loop; namely, the slices are finite. That implies the following two: first, the condition for \Cref{claim:decodestopnoise} is valid, and therefore the replacement argument can be used to infer a correction for logical measurement gadgets. Second, the errors are finite, and the sweep rule is, in some sense, effective (\Cref{claim:polesmove}), and soon we will see that it sets a structure for the histories of slices.


That motivates us to describe the evolution of a fault to a given step and define the notion of a step that admits infinite error.

\begin{definition}[Fault propagation to a step $\tau$]
  \label{def:localprop}
  Let $C=\{u_i\}_i$ be a quantum circuit, let $J=(j_1,\ldots,j_k)$ be an order agreeing with $\pi_C$, and let $u_i\in C$. Suppose that $u_i=u_{j_l}$ in this order. If $l>\tau$, we define the propagation of $u_i$ to step $\tau$ to be the identity operator.
Otherwise, define operators $\mathcal{D}_0,\mathcal{D}_1,\ldots,\mathcal{D}_{\tau-l}$ recursively by setting $\mathcal{D}_0=u_{j_l}$ and requiring that for every $1\le s\le \tau-l$,
  \begin{equation*}
    \begin{split}
      \mathcal{D}_{s-1}u_{j_{l+s}} = u_{j_{l+s}}\mathcal{D}_s.
    \end{split}
  \end{equation*}
  We then define the \emph{propagation of $u_i$ to step $\tau$} to be the operator
  \begin{equation*}
    \begin{split}
      \mathcal{D} \colonequals \mathcal{D}_{\tau-l}.
    \end{split}
  \end{equation*}
Now let $C[\mathcal{E}]$ be a noisy circuit, where all of its quantum registers are encoded in the $(d,d^{\prime})$-Toric code with side length $n$. Let $J$ be an order agreeing with $\pi_{C[\mathcal{E}]}$, and let $\tau$ be a step in that order. For each faulty gate $e \in \mathcal{E}$ that occurs no later than step $\tau$, denote by $\mathcal{D}^{\tau}_{e}$ its propagation to step $\tau$. We define the propagated error at step $\tau$ by multiplying these propagated faults in reverse topological order.
  \begin{equation*}
    \begin{split}
      \mathcal{D}^{\tau} \colonequals \prod_{e\preceq \tau} \mathcal{D}^{\tau}_{e}.
    \end{split}
  \end{equation*}
Moreover, we say that a step $\tau$ admits an infinite error if there is a register such that, in the view of the $\infty$-extension of $C[\mathcal{E}]$, the Pauli decomposition of the propagated error at step $\tau$ has a term that is infinite (\Cref{def:finiteerr}).
\end{definition}


\begin{proof}
Let $\mathcal{H}_{\tau}$ and $\mathcal{H}_{\tau+1}$ be the analysis Toom's contours at steps $\tau$ and $\tau+1$. And let $P_{1}, \ldots, P_{t}, \tilde{P}_{1}$ and $P_{1}, \ldots, P_{t}, \tilde{P}_{2}$ be the reduced Pauli-paths (up to a step) that induce Toom's contours. In addition, let $g$ be the gate that is applied at step $\tau+1$. Assume that step $\tau+1$ is the first to admit an infinite slice. For the sake of the proof, let us denote by $C[\mathcal{E}]|_{\rightarrow \tau}$ the quantum circuit up to step $\tau$, so
  \begin{equation*}
    \begin{split}
 C[\mathcal{E}]|_{\rightarrow \tau+1} = g C[\mathcal{E}]|_{\rightarrow \tau}  
    \end{split}
  \end{equation*}
By definition of Pauli-path (\Cref{def:paulipath}), we have that:
\begin{equation*}
  \begin{split}
    \Trr{} \big( (C[\mathcal{E}]|_{\tau} \circ P_{t}) \tilde{P}_{1}   \big) \neq 0 
  \end{split}
\end{equation*}
Thus, there are constants $\gamma, \bar{\gamma} \in \mathbb{R}$ and an operator $\mathbf{A} \in \mathcal{L}$, such that $\Trr{} (\mathbf{A} \tilde{P}_{1}) = 0$ and:
\begin{equation*}
  \begin{split}
  C[\mathcal{E}]|_{\tau} \circ P_{t} = \gamma \tilde{P}_{1} + \bar{\gamma} \mathbf{A}
  \end{split}
\end{equation*}
\inb{\textbf{Think what to write here.}}
\begin{equation*}
  \begin{split}
    &\Trr{} \big( (C[\mathcal{E}]|_{\tau + 1} \circ P_{t}) \tilde{P}_{2}   \big) = \Trr{} \big(   (gC[\mathcal{E}]|_{\tau}  \circ P_{t}) \tilde{P}_{2}   \big) = \gamma \Trr{} \big( (g \circ \tilde{P}_{1} )  \tilde{P}_{2} \big) +  \bar{\gamma} \Trr{} \big (g \circ \mathbf{A} )  \tilde{P}_{2}   \big) 
  \end{split}
\end{equation*}
We argue that we can assume, without loss of generality, that the left term is non-zero. To see this, suppose that it vanishes; then there has to be another Pauli term $\tilde{P} \in \mathcal{P}$ in the decomposition of $\mathbf{A}$ for which:
\begin{equation*}
    \begin{split}
      \Trr{} \big( (g \circ \tilde{P} )  \tilde{P}_{2} \big) \neq 0
  \end{split}
\end{equation*}
Otherwise, the whole trace is zeroing out in contradiction to the fact that $P_{1},..,P_{t},\tilde{P_{2}}$ is a Pauli-path. So, we could initially choose the Pauli-path that induces '$\mathcal{H}_{\tau}$'\footnote{Observe that the 'reinduced' Toom's counter is different from $\mathcal{H}_{\tau}$.} to be $P_{1},..,P_{t},\tilde{P}$.



Now, let $K_{\infty}$ be a slice in $\mathcal{H}_{\tau+1}$ that is not in $\mathcal{H}_{\tau+1}$, and let $K$ be a finite subset of $K_{\infty}$ with size exactly $n$. Observe that $K$ can be chosen to support $g$'s locations. The vertices in $K$ support the faces of $\tilde{P_{2}} = g \circ \tilde{P}$. However, since $g$ acts over at most $\Delta$ faces, there are at most $\Delta$ faces that support $\tilde{P}_{2}$ and don't support $\tilde{P}_{1}$. Furthermore, since the arrow connection rules define a connection of less than $d^{2}$ arrows per face, the arrows in $\mathcal{H}_{\tau+1}$ can be obtained by the addition of at most $\Delta d^{2}$ arrows to the arrows of $\mathcal{H}_{\tau}$. Thus, any connected (by arrows) subset of vertices in $\mathcal{H}_{t+1}$ can be covered by the merging of at most $\Delta d^{2}+1$ connected components in $\mathcal{H}_{\tau}$ with the vertices supporting $g$ .


To finish, imagine the space-time graph, consisting of three time-layers. In the bottom layer, set delegates vertices, each corresponding to a slice in the last layer of $\mathcal{H}_{\tau}$. In the second layer, set the vertices in those slices and the vertices supporting $g$'s locations. Connect each of the vertices of $\mathcal{H}_{\tau}$'s slices to its delegate with an arrow.
And in the final layer, set the vertices of $K$, connected to the vertices in the second layer by either connecting the vertex to itself \inb{\textbf{(Change it)}} or by the additional arrows. Finally, add the required forks so that the result is indeed a space-time graph.

Denote by $K^{\prime}$ the slice in the last layer of $\mathcal{H}_{\tau}$ that is maximal in size. Using the \Cref{claim:prevslices}, combining that the vertices of the third layer contained in the second gives:  
\begin{equation*}
  \begin{split}
     & \Size{K}  \le \Delta d^{2} ( \Size{K^{\prime}} +  \Size{\text{fork}} ) \\
     \Rightarrow & \Size{K^{\prime}} \le \frac{1}{\Delta d^{2}} \Size{K} - \Size{ \text{fork} } \le \alpha \Size{K} 
  \end{split}
\end{equation*}


\end{proof}

\inb{\textbf{We have one crucial problem: the Toom's contour is defined over the $\infty$-extension, so there are an infinite number of $g$'s.}}


\begin{claim}
  \label{claim:intersyn}
  Let $\alpha = (\Delta d^{2})^{-1}$. Denote by $A$ the event that there is a step that admits an infinite slice. In addition, for any step $\tau \in J$, denote by $B_{\tau}$ the event that step $\tau$ is the first step to addmit a slice at size $\alpha n$. Denote by $w$ and $\mu$ the width and the depth of $C[\mathcal{E}]$. Then the probability that $A$ occurs is bounded by:
  \begin{equation*}
    \begin{split}
      \prb{A} \le w\mu\max_{\tau\in[w\mu]}\prb{B_{\tau}}.
    \end{split}
  \end{equation*}
\end{claim}

\begin{proof}
  For the sake of the proof, we say that a chain $x$ is in operator $O$ if, when considering the Pauli decomposition of the image of $O$ in the $\infty$-extension, there is a correspondence between the faces of $x$ and the (partial) support of one of the terms in $O$'s image.

  Denote by $x$ the maximal, in size, reduced connected $d^{\prime}$-dimensional chain in $\mathcal{D}^{\tau+1}$. Assume that $\tau$ is the first step to admits an infinite error, thus any chain in $\mathcal{D}^{\tau}$ has a size lower than $n$. Denote by $u$ the gates which is applied at step $\tau$, so there have to be Paulis $P_{1},P_{2}$ such that: 
  \begin{equation*}
    \begin{split}
      P_{1} \in \text{Supp } \mathcal{D}^{\tau}, \, \,  P_{2} \in \text{Supp } \mathcal{D}^{\tau + 1}, \, \, x \text{ is both in } P_{1},P_{2}  \, \, \text{ and } \expval*{P_{1}|u|P_{2}} \neq 0
    \end{split}
  \end{equation*}
  Since $u$ has a fan at most $\Delta$, and since any qubit ($d^{\prime}$-dimensional face) is connected to at most $d^{2}$ qubits (by stepping through a check\footnote{a $d^{\prime}\pm1$ dimensional face}), we get that the application of $u$ on $P_{1}$ touches at most $\alpha^{-1}$ chains, and the result of the modification yields $x$. Therefore, $x$ resides inside a union of at most $\alpha^{-1}$, and since $x$ is connected, we have that its size is bounded by summing the size of each chain in the union (\Cref{claim:tringinq}). Therefore, if any of the chains is at size at most $\alpha n$, it implies that the size of $x$ is lower than $n$. Hence, a necessary condition to step $\tau+1$ to be the first step to admit an infinite error is that step $\tau-1$ admits errors at size greater than $\alpha n$, and hence the existence of a step $\tau^{\prime}\le \tau$ such that $B_{\tau^{\prime}}$ occurs. The inequality follows by using the union bound over all the steps in $C[\mathcal{E}]$.
\end{proof}



